\documentclass[11pt]{article}
\usepackage[utf8]{inputenc}
\usepackage[T1]{fontenc}
\usepackage[spanish]{babel}
\usepackage{amsmath, amssymb, amsthm, mathtools}
\usepackage{braket}
\usepackage{dsfont}
\usepackage{graphicx}
\usepackage{float}
\usepackage[margin=2.5cm]{geometry}
\usepackage{hyperref}
\hypersetup{hidelinks}
\emergencystretch=2em
\allowdisplaybreaks

\newcommand{\ii}{\mathrm{i}}
\newcommand{\ee}{\mathrm{e}}
\newcommand{\Id}{\mathds{1}}

\title{Estados coherentes fot\'onico y fon\'onico}
\author{Ricardo Del Rosal Gardu\~no \\ \small Doctorado en F\'isica, BUAP}
\date{Notas de doctorado}

\begin{document}
\maketitle
\tableofcontents
\bigskip

\section{Estado inicial}

Se considera el modelo de Holstein--Cummings despu\'es del desplazamiento
vibracional, de la transformaci\'on a la base polarit\'onica y de la
aproximaci\'on de onda rotante que conduce al hamiltoniano efectivo
$\hat H_{\mathrm{ef}}^{(n)}$. El estado inicial se toma como un producto entre
un estado coherente fot\'onico, un estado electr\'onico molecular general y un
estado coherente fon\'onico:
\begin{equation}
  \boxed{
  \ket{\tilde\Psi(0)}
  =
  \ket{\alpha}
  \otimes
  (c_e\ket e+c_g\ket g)
  \otimes
  \ket{\tilde\beta}.
  }
\end{equation}
La tilde en $\ket{\tilde\Psi}$ y en $\ket{\tilde\beta}$ indica que se est\'a
trabajando en la representaci\'on desplazada del modo vibracional. Los estados
coherentes se expanden como
\begin{equation}
  \boxed{
  \ket{\alpha}
  =
  \sum_{n=0}^{\infty}c_n\ket n,
  \qquad
  c_n
  =
  \ee^{-|\alpha|^2/2}
  \frac{\alpha^n}{\sqrt{n!}},
  }
\end{equation}
y
\begin{equation}
  \boxed{
  \ket{\tilde\beta}
  =
  \sum_{m=0}^{\infty}\tilde c_m\ket m,
  \qquad
  \tilde c_m
  =
  \ee^{-|\tilde\beta|^2/2}
  \frac{\tilde\beta^m}{\sqrt{m!}}.
  }
\end{equation}
Los coeficientes $c_n$ pertenecen al modo fot\'onico de la cavidad, mientras
que $\tilde c_m$ pertenecen al modo fon\'onico desplazado. Esta distinci\'on es
importante porque en las expresiones finales aparecen productos de ambos tipos
de amplitudes.

\section{Base polarit\'onica}

En cada bloque de excitaci\'on fija $n\geq1$, la base desnuda
$\{\ket{e,n-1},\ket{g,n}\}$ se reescribe en la base polarit\'onica
\begin{equation}
  \boxed{
  \ket{+,n}
  =
  \frac{1}{\sqrt2}(\ket{e,n-1}+\ket{g,n}),
  \qquad
  \ket{-,n}
  =
  \frac{1}{\sqrt2}(\ket{e,n-1}-\ket{g,n}).
  }
\end{equation}
Por tanto,
\begin{equation}
  \ket{e,n-1}
  =
  \frac{\ket{+,n}+\ket{-,n}}{\sqrt2},
  \qquad
  \ket{g,n}
  =
  \frac{\ket{+,n}-\ket{-,n}}{\sqrt2}.
\end{equation}
Al proyectar la parte opto-electr\'onica inicial sobre los polaritones se
obtienen los coeficientes
\begin{equation}
  \boxed{
  d_n^\pm
  =
  \frac{1}{\sqrt2}
  (c_{n-1}c_e\pm c_nc_g).
  }
\end{equation}
El signo superior mide la amplitud inicial en la rama polarit\'onica
$\ket{+,n}$ y el signo inferior la amplitud inicial en la rama
$\ket{-,n}$. La dependencia con $c_{n-1}$ y $c_n$ aparece porque el estado
$\ket{e,n-1}$ tiene un fot\'on menos que el estado $\ket{g,n}$ dentro del mismo
bloque de excitaci\'on total $n$.

\section{Parte din\'amica del estado inicial}

El hamiltoniano efectivo polarit\'on--fon\'on conecta los pares
\begin{equation}
  \boxed{
  \ket{+,n,m-1}
  \longleftrightarrow
  \ket{-,n,m},
  \qquad n\geq1,\quad m\geq1.
  }
\end{equation}
Por ello, la parte din\'amica del estado inicial se ordena agrupando amplitudes
que pertenecen al mismo doblete:
\begin{equation}
  \boxed{
  \begin{aligned}
  \ket{\tilde\Psi_{\mathrm{dyn}}(0)}
  &=
  \sum_{n=1}^{\infty}
  \sum_{m=1}^{\infty}
  \left(
  \tilde c_{m-1}d_n^+\ket{+,n,m-1}
  +
  \tilde c_md_n^-\ket{-,n,m}
  \right)
  \\
  &\quad+
  \sum_{n=1}^{\infty}
  \tilde c_0d_n^-\ket{-,n,0}.
  \end{aligned}
  }
\end{equation}
El segundo t\'ermino contiene los estados $\ket{-,n,0}$. Estos estados no tienen
pareja $\ket{+,n,-1}$, por lo que quedan desacoplados dentro de la din\'amica
efectiva. Tambi\'en existe el sector $c_0c_g\ket{g,0}\ket{\tilde\beta}$,
asociado con cero excitaciones opto-electr\'onicas. Ese sector queda separado
de los dobletes $n\geq1$ y solamente contribuye a la inversi\'on electr\'onica
como poblaci\'on del estado base.

\section{Evoluci\'on temporal}

La evoluci\'on de la parte din\'amica conserva la estructura por dobletes:
\begin{equation}
  \boxed{
  \begin{aligned}
  \ket{\tilde\Psi_{\mathrm{dyn}}(t)}
  &=
  \sum_{n=1}^{\infty}
  \sum_{m=1}^{\infty}
  \left(
  p_{nm}^{+}(t)\ket{+,n,m-1}
  +
  p_{nm}^{-}(t)\ket{-,n,m}
  \right)
  \\
  &\quad+
  \sum_{n=1}^{\infty}
  p_{n0}^{-}(t)\ket{-,n,0}.
  \end{aligned}
  }
\end{equation}
Para $m\geq1$, las amplitudes son
\begin{equation}
  \boxed{
  \begin{aligned}
  p_{nm}^{+}(t)
  &=
  \ee^{-\ii A_{nm}t}
  \Bigg[
  \left(
  \cos\frac{\Omega_{nm}^{\mathrm{ef}}t}{2}
  -
  \ii\frac{\delta_n}{\Omega_{nm}^{\mathrm{ef}}}
  \sin\frac{\Omega_{nm}^{\mathrm{ef}}t}{2}
  \right)
  \tilde c_{m-1}d_n^+
  \\
  &\hspace{3.0cm}
  +
  \ii\frac{\kappa_m}{\Omega_{nm}^{\mathrm{ef}}}
  \sin\frac{\Omega_{nm}^{\mathrm{ef}}t}{2}
  \tilde c_md_n^-
  \Bigg],
  \\
  p_{nm}^{-}(t)
  &=
  \ee^{-\ii A_{nm}t}
  \Bigg[
  \ii\frac{\kappa_m}{\Omega_{nm}^{\mathrm{ef}}}
  \sin\frac{\Omega_{nm}^{\mathrm{ef}}t}{2}
  \tilde c_{m-1}d_n^+
  \\
  &\hspace{3.0cm}
  +
  \left(
  \cos\frac{\Omega_{nm}^{\mathrm{ef}}t}{2}
  +
  \ii\frac{\delta_n}{\Omega_{nm}^{\mathrm{ef}}}
  \sin\frac{\Omega_{nm}^{\mathrm{ef}}t}{2}
  \right)
  \tilde c_md_n^-
  \Bigg].
  \end{aligned}
  }
\end{equation}
Las cantidades que aparecen en estas amplitudes son
\begin{equation}
  \boxed{
  A_{nm}
  =
  n\Omega
  +
  \frac{\nu\lambda^2}{4}
  +
  \nu\left(m-\frac12\right),
  \qquad
  \delta_n=2g\sqrt n-\nu,
  }
\end{equation}
\begin{equation}
  \boxed{
  \kappa_m=\nu\lambda\sqrt m,
  \qquad
  \Omega_{nm}^{\mathrm{ef}}
  =
  \sqrt{\delta_n^2+\kappa_m^2}.
  }
\end{equation}
El t\'ermino desacoplado evoluciona solamente con su fase propia:
\begin{equation}
  \boxed{
  p_{n0}^{-}(t)
  =
  \tilde c_0d_n^-
  \exp\left[
  -\ii
  \left(
  n\Omega+\frac{\nu\lambda^2}{4}-g\sqrt n
  \right)t
  \right].
  }
\end{equation}

\section{Inversi\'on polarit\'onica}

La inversi\'on polarit\'onica efectiva mide la diferencia entre poblaci\'on de
la rama superior y poblaci\'on de la rama inferior:
\begin{equation}
  \mathcal W(t)
  =
  \bra{\tilde\Psi(t)}
  \hat S_z
  \ket{\tilde\Psi(t)}.
\end{equation}
En la base polarit\'onica,
\begin{equation}
  \hat S_z^{(n)}
  =
  \ket{+,n}\bra{+,n}
  -
  \ket{-,n}\bra{-,n}.
\end{equation}
Por tanto, cada amplitud $p_{nm}^{+}$ entra con signo positivo y cada amplitud
$p_{nm}^{-}$ con signo negativo. El estado desacoplado $\ket{-,n,0}$ tambi\'en
pertenece a la rama inferior, de modo que
\begin{equation}
  \boxed{
  \mathcal W(t)
  =
  \sum_{n=1}^{\infty}
  \sum_{m=1}^{\infty}
  \left(
  |p_{nm}^{+}(t)|^2
  -
  |p_{nm}^{-}(t)|^2
  \right)
  -
  \sum_{n=1}^{\infty}
  |p_{n0}^{-}(t)|^2.
  }
\end{equation}
Esta expresi\'on no contiene coherencias entre $+$ y $-$ porque
$\hat S_z$ es diagonal en la base polarit\'onica.

\section{Inversi\'on electr\'onica}

La inversi\'on electr\'onica se define con el observable
\begin{equation}
  \hat\sigma_z
  =
  \ket e\bra e-\ket g\bra g,
  \qquad
  W(t)
  =
  \bra{\tilde\Psi(t)}
  \hat\sigma_z
  \ket{\tilde\Psi(t)}.
\end{equation}
En la base polarit\'onica, $\hat\sigma_z$ no es diagonal. Para un mismo bloque
$n$ intercambia las ramas $\ket{+,n}$ y $\ket{-,n}$ sin cambiar el n\'umero
fon\'onico. Por esta raz\'on deben emparejarse amplitudes que tengan el mismo
n\'umero fon\'onico $m$.

La amplitud $p_{nm}^{-}$ multiplica a $\ket{-,n,m}$, mientras que
$p_{n,m+1}^{+}$ multiplica a $\ket{+,n,m}$. Por tanto, la inversi\'on
electr\'onica de los dobletes $n\geq1$ se obtiene de los productos
$p_{n,m+1}^{+*}p_{nm}^{-}$. A esto se suma el sector desacoplado
$c_0c_g\ket{g,0}\ket{\tilde\beta}$, que contribuye con signo negativo porque
est\'a en el estado electr\'onico base. As\'i,
\begin{equation}
  \boxed{
  W(t)
  =
  2\operatorname{Re}
  \left[
  \sum_{n=1}^{\infty}
  \sum_{m=0}^{\infty}
  p_{n,m+1}^{+*}(t)p_{nm}^{-}(t)
  \right]
  -
  |c_0c_g|^2.
  }
\end{equation}
El detalle esencial es que $\hat\sigma_z$ no cambia $m$. Si se emparejaran
amplitudes con n\'umeros fon\'onicos distintos, el producto interno se anular\'ia
por ortogonalidad de los estados de Fock.

\section{Casos iniciales y par\'ametros num\'ericos}

En las gr\'aficas se toman estados coherentes en ambos modos con
\begin{equation}
  \boxed{
  \alpha=\beta=3.
  }
\end{equation}
Como la derivaci\'on est\'a escrita en la representaci\'on desplazada, la
amplitud que aparece en la expansi\'on fon\'onica se denota por
$\tilde\beta$. En esta secci\'on se usa el mismo valor num\'erico,
$\tilde\beta=3$.
Por tanto, las distribuciones fot\'onica y fon\'onica est\'an centradas en
$|\alpha|^2=|\tilde\beta|^2=9$. El estado electr\'onico inicial se estudia en
dos preparaciones:
\begin{equation}
  \boxed{
  \text{E:}\quad c_e=1,\qquad c_g=0,
  }
\end{equation}
y
\begin{equation}
  \boxed{
  \text{S:}\quad c_e=-c_g=\frac{1}{\sqrt2}.
  }
\end{equation}
En el caso E, el sistema parte en el estado electr\'onico excitado. En el caso
S, el estado opto-electr\'onico se proyecta, cerca del m\'aximo de la
distribuci\'on coherente, principalmente sobre una rama polarit\'onica.

Los par\'ametros usados son
\begin{equation}
  \boxed{
  g=1,
  \qquad
  \Omega=1000g,
  \qquad
  \nu=6g,
  \qquad
  \eta=\nu\lambda=0.02g.
  }
\end{equation}
La elecci\'on $\nu=6g$ corresponde a la resonancia polarit\'on--fon\'on del
bloque central $n=9$, porque
\begin{equation}
  \delta_9=2g\sqrt9-\nu=6g-\nu=0.
\end{equation}
Para una distribuci\'on fon\'onica centrada en $M=9$, el tiempo de revival
obtenido con la aproximaci\'on coherente fon\'onica es
\begin{equation}
  \boxed{
  t_r
  =
  \frac{4\pi\sqrt M}{\nu\lambda}
  =
  \frac{12\pi}{0.02}
  \simeq
  1884.96.
  }
\end{equation}

\section{Resultados gr\'aficos}

\subsection{Inversi\'on polarit\'onica}

\begin{figure}[H]
  \centering
  \begin{minipage}{0.49\linewidth}
    \centering
    \includegraphics[width=\linewidth]{\detokenize{../../notas-doctorado/imgs/Estado coherente fotonico y fononico/InversionP_E.png}}
    \small (a) Estado E.
  \end{minipage}
  \hfill
  \begin{minipage}{0.49\linewidth}
    \centering
    \includegraphics[width=\linewidth]{\detokenize{../../notas-doctorado/imgs/Estado coherente fotonico y fononico/InversionP_S.png}}
    \small (b) Estado S.
  \end{minipage}
  \caption{Inversi\'on polarit\'onica para estados iniciales fot\'onico y
  fon\'onico coherentes con $\alpha=\beta=3$. En la representaci\'on desplazada
  se escribe $\tilde\beta=3$. En (a) se usa
  $c_e=1$, $c_g=0$; en (b), $c_e=-c_g=1/\sqrt2$. Ambas se\~nales muestran el
  comportamiento esperado de colapsos, revivals y ensanchamiento de los
  l\'obulos. La diferencia es que el caso S sigue con mayor precisi\'on el
  tiempo anal\'itico $t_r=1884.96$. La raz\'on es que, cerca del m\'aximo de la
  distribuci\'on fot\'onica, $c_{n-1}\simeq c_n$; entonces
  $d_n^+=(c_{n-1}-c_n)/2$ queda suprimido y domina
  $d_n^-=(c_{n-1}+c_n)/2$. La preparaci\'on se parece m\'as al caso
  polarit\'onico ideal usado para estimar el tiempo de revival.}
\end{figure}

\subsection{Inversi\'on electr\'onica}

\begin{figure}[H]
  \centering
  \begin{minipage}{0.49\linewidth}
    \centering
    \includegraphics[width=\linewidth]{\detokenize{../../notas-doctorado/imgs/Estado coherente fotonico y fononico/InversionE_E.png}}
    \small (a) Estado E.
  \end{minipage}
  \hfill
  \begin{minipage}{0.49\linewidth}
    \centering
    \includegraphics[width=\linewidth]{\detokenize{../../notas-doctorado/imgs/Estado coherente fotonico y fononico/InversionE_S.png}}
    \small (b) Estado S.
  \end{minipage}
  \caption{Inversi\'on electr\'onica para las mismas preparaciones iniciales. En
  el caso E, la inversi\'on comienza cerca de $1$ porque la mol\'ecula parte en
  $\ket e$. La din\'amica reparte la poblaci\'on entre car\'acter excitado y base,
  y las muchas frecuencias asociadas con las distribuciones coherentes producen
  colapsos y revivals visibles alrededor de cero. En el caso S, las poblaciones
  electr\'onicas iniciales est\'an balanceadas y la proyecci\'on dominante cerca de
  $n=9$ corresponde a una rama polarit\'onica con mezcla casi equilibrada de
  $\ket e$ y $\ket g$. Por eso la inversi\'on queda suprimida y oscila con
  amplitud mucho menor. Las oscilaciones electr\'onicas son m\'as r\'apidas que
  las polarit\'onicas porque dependen de coherencias entre ramas con el mismo
  n\'umero fon\'onico, donde aparecen fases de escala $\nu$, mientras que la
  envolvente polarit\'onica est\'a gobernada por escalas efectivas del orden de
  $\eta\sqrt m$.}
\end{figure}

\subsection{Comparaci\'on con el hamiltoniano completo}

\begin{figure}[H]
  \centering
  \begin{minipage}{0.49\linewidth}
    \centering
    \includegraphics[width=\linewidth]{\detokenize{../../notas-doctorado/imgs/Estado coherente fotonico y fononico/ComparacionPE.png}}
    \small (a) $\mathcal W(t)$, estado E.
  \end{minipage}
  \hfill
  \begin{minipage}{0.49\linewidth}
    \centering
    \includegraphics[width=\linewidth]{\detokenize{../../notas-doctorado/imgs/Estado coherente fotonico y fononico/ComparacionPS.png}}
    \small (b) $\mathcal W(t)$, estado S.
  \end{minipage}

  \vspace{0.4cm}

  \begin{minipage}{0.49\linewidth}
    \centering
    \includegraphics[width=\linewidth]{\detokenize{../../notas-doctorado/imgs/Estado coherente fotonico y fononico/ComparacionEE.png}}
    \small (c) $W(t)$, estado E.
  \end{minipage}
  \hfill
  \begin{minipage}{0.49\linewidth}
    \centering
    \includegraphics[width=\linewidth]{\detokenize{../../notas-doctorado/imgs/Estado coherente fotonico y fononico/ComparacionES.png}}
    \small (d) $W(t)$, estado S.
  \end{minipage}
  \caption{Comparaci\'on entre la evoluci\'on anal\'itica generada por el
  hamiltoniano efectivo y la evoluci\'on num\'erica del hamiltoniano completo
  $H_{HC}$. Se muestran la inversi\'on polarit\'onica para los estados E y S, y
  la inversi\'on electr\'onica para las mismas dos preparaciones. En los cuatro
  casos la coincidencia es alta durante el intervalo mostrado, lo que valida el
  uso de $\hat H_{\mathrm{ef}}^{(n)}$ para describir la din\'amica relevante en
  este r\'egimen resonante.}
\end{figure}

\subsection{Fidelidad}

\begin{figure}[H]
  \centering
  \begin{minipage}{0.49\linewidth}
    \centering
    \includegraphics[width=\linewidth]{\detokenize{../../notas-doctorado/imgs/Estado coherente fotonico y fononico/FidelidadE.png}}
    \small (a) Estado E.
  \end{minipage}
  \hfill
  \begin{minipage}{0.49\linewidth}
    \centering
    \includegraphics[width=\linewidth]{\detokenize{../../notas-doctorado/imgs/Estado coherente fotonico y fononico/FidelidadS.png}}
    \small (b) Estado S.
  \end{minipage}
  \caption{Fidelidad entre el estado evolucionado con el hamiltoniano efectivo
  y el estado evolucionado num\'ericamente con el hamiltoniano completo. En el
  caso E, la preparaci\'on inicial excita de manera comparable las ramas
  polarit\'onicas superior e inferior y, por la distribuci\'on coherente
  fot\'onica, involucra muchos bloques $n$. Las peque\~nas fases y acoplamientos
  descartados por la aproximaci\'on efectiva se acumulan con mayor rapidez. En
  el caso S, cerca del m\'aximo de la distribuci\'on domina una rama
  polarit\'onica y la evoluci\'on queda m\'as cerca del r\'egimen resonante ideal;
  por eso la fidelidad se mantiene alta durante m\'as tiempo.}
\end{figure}

\section{Lectura f\'isica final}

Esta secci\'on generaliza el an\'alisis de un bloque polarit\'onico ideal al caso
en que tanto el campo fot\'onico como el modo fon\'onico tienen distribuciones
coherentes. La din\'amica ya no est\'a centrada en un solo par $(n,m)$, sino en
una ventana de bloques alrededor de $n\simeq|\alpha|^2$ y
$m\simeq|\tilde\beta|^2$. Cada bloque aporta una frecuencia efectiva
$\Omega_{nm}^{\mathrm{ef}}$, y la suma de todas ellas produce colapsos,
revivals y envolventes anchas.

La preparaci\'on electr\'onica controla cu\'anto peso inicial cae en cada rama
polarit\'onica. Esa proyecci\'on explica por qu\'e el caso S sigue mejor el
tiempo de revival calculado para una rama polarit\'onica dominante, mientras que
el caso E muestra una mezcla m\'as amplia de contribuciones. La comparaci\'on
con el hamiltoniano completo confirma que, en el r\'egimen usado, el
hamiltoniano efectivo conserva la estructura din\'amica esencial.

\end{document}
